%!TeX program = lualatex
\renewcommand{\thelektion}{L05 \textperiodcentered{} Bilder \& Farben}

\begin{frame}\lessontitle{Lektion 05}{Kodierung von Bildern \& Farben}{Pixel, RGB und Farbtiefe \textperiodcentered{} Skript, Kapitel 2}\end{frame}

\begin{frame}[fragile]\FDUtitle{Ein Schwarz-Weiss-Bild}
\begin{table}[H]
	\centering
	\begin{tabular}{lllll}
		1 & 0 & 0 & 0 & 1 \\
		0 & 1 & 0 & 1 & 0 \\
		0 & 0 & 1 & 0 & 0 \\
		0 & 1 & 0 & 1 & 0 \\
		1 & 0 & 0 & 0 & 1
	\end{tabular}
\end{table}
\begin{center}
	Welches Bild wird dargestellt?\\
	Wie viel Bit Speicherplatz benötigt es?
\end{center}
\only<2->{\begin{center}\color{FDUteal}\bfseries Ein Kreuz \textperiodcentered{} $5\cdot 5 = 25$ Bit\end{center}}
\end{frame}

\begin{frame}\FDUtitle{Vom Pixel zum Bildschirm}
\begin{columns}[c,onlytextwidth]
\begin{column}{.46\textwidth}
\centering
\includegraphics[width=\linewidth]{Grundlagen_Info/13_ZahlendarstellungenUndKodierungen/Figures/pixels}
\end{column}
\begin{column}{.46\textwidth}
\tib{Pixel} = \enquote{picture} + \enquote{element}
\vspace{2mm}
\begin{itemize}\setlength\itemsep{5pt}
  \item Winzige rot/grün/blau leuchtende \glspl{LED}
  \item Kleinstes Bildelement eines Bildschirms
\end{itemize}
\end{column}
\end{columns}
\end{frame}

\begin{frame}\FDUtitle{Additive Farbmischung}
\begin{figure}[H]
	\centering
	\begin{subfigure}{.49\textwidth}
		\includegraphics[width=\textwidth]{Grundlagen_Info/13_ZahlendarstellungenUndKodierungen/Figures/rgb-farbmischung-1}
	\end{subfigure}
	\begin{subfigure}{.49\textwidth}
		\includegraphics[width=\textwidth]{Grundlagen_Info/13_ZahlendarstellungenUndKodierungen/Figures/rgb-farbmischung-2}
	\end{subfigure}
\end{figure}
\begin{center}Jede Farbe wird durch drei Werte (Rot, Grün, Blau) zwischen 0 und 255 kodiert.\end{center}
\end{frame}

\begin{frame}[fragile]\FDUtitle{RGB und Hexadezimal}
24 Bit pro Farbe (je 8 Bit für R, G, B) $\rightarrow$ hexadezimal viel kompakter:
\vspace{3mm}

\texttt{\textcolor{red}{11111111}\textcolor{green!60!black}{00000000}\textcolor{blue}{00000000} = \#\textcolor{red}{FF}\textcolor{green!60!black}{00}\textcolor{blue}{00} = \textcolor{red}{255}, \textcolor{green!60!black}{0}, \textcolor{blue}{0}} = reines Rot

\texttt{\textcolor{red}{10000000}\textcolor{green!60!black}{10000000}\textcolor{blue}{10000000} = \#\textcolor{red}{80}\textcolor{green!60!black}{80}\textcolor{blue}{80} = \textcolor{red}{128}, \textcolor{green!60!black}{128}, \textcolor{blue}{128}} = grau
\vspace{3mm}
\begin{center}
Wie viele Farben lassen sich damit insgesamt darstellen? $\quad 2^{24}\approx 16.7$ Mio.
\end{center}
\end{frame}

\begin{frame}\FDUtitle{Farbtiefe}
\begin{figure}[H]
	\centering
	\includegraphics[width=.85\linewidth]{Grundlagen_Info/13_ZahlendarstellungenUndKodierungen/Figures/colordepth}
\end{figure}
\end{frame}

\begin{frame}\FDUtitle{Farbtiefe: Beispiele}
\begin{figure}[H]
	\centering
	\begin{subfigure}{.48\textwidth}
		\centering
		\includegraphics[height=38mm,keepaspectratio]{Grundlagen_Info/13_ZahlendarstellungenUndKodierungen/Figures/colordepth-bar}
	\end{subfigure}\hfill
	\begin{subfigure}{.48\textwidth}
		\centering
		\includegraphics[height=38mm,keepaspectratio]{Grundlagen_Info/13_ZahlendarstellungenUndKodierungen/Figures/colordepth-cat}
	\end{subfigure}
\end{figure}
\vspace{-2mm}
\begin{center}\small Je geringer die Farbtiefe, desto weniger fein abgestuft.\end{center}
\end{frame}

\begin{frame}\FDUtitle{GIF und Farbpaletten}
\begin{center}
\gls{GIF}-Bilder unterscheiden meist 256 Farben. Wie viele Bits pro Pixel braucht das?
\end{center}
\only<2->{
\begin{center}\Large 8 Bit\end{center}
\begin{center}Statt für jedes Pixel dieselbe Anzahl Bits pro Farbe zu speichern, wird oft ein \tib{Index} auf eine Farbpalette verwendet.\end{center}}
\end{frame}

\begin{frame}{Auftrag}{Kapitel 2: Bilder \& Farben}
\aufgabebox{\textbf{Übungen}: Schwarz-Weiss-Bild, Pixel-Speicherbedarf \& GIF-Palette}
\vspace{3mm}
\aufgabebox{\textbf{Übungen}: Hexadezimale Farbzuordnung \& Speicherbedarf $500\times 500$ Bild (64 Farben)}
\vspace{3mm}
\challengebox{Weitere Farben erraten}
\end{frame}
